\section{El grafo de decisiones}
\label{sec:grafo}

\subsection{Forma}

La campaña se representa como un grafo de diez nodos en diez etapas ordenadas,
cada uno con su pregunta y su palabra de firmeza. No es un diagrama ilustrativo:
es un objeto que se calcula a partir del registro, y cada nodo publica junto a su
palabra el valor en el que está parado, campo por campo. Una palabra dice con qué
firmeza se sostiene una decisión y no dice nada sobre \emph{qué} se decidió; un
lector que no vea el valor no puede distinguir un valor por omisión heredado de
una elección deliberada que resulta no estar medida. Las dos leen
\cod{inherited} y solo una de las dos es una sorpresa.

\begin{center}
\begin{tabularx}{\linewidth}{@{}r L{2.6cm} X@{}}
\toprule
etapa & nodo & pregunta\\
\midrule
0 & \cod{frame}       & En qué ventana, pivote y régimen de accretion se lee cada número\\
1 & \cod{substrate}   & En qué representación se calcula el vecindario\\
2 & \cod{bank}        & De qué corpus vienen los donantes, y bajo qué política\\
3 & \cod{retriever}   & Cómo se extraen candidatos del banco, y a qué profundidad\\
4 & \cod{generator}   & Si algún candidato llega sin donante\\
5 & \cod{scoring}     & Qué ponderación convierte un candidato en una puntuación\\
6 & \cod{features}    & Qué rasgos por candidato entran en un modelo\\
7 & \cod{reranking}   & Si un modelo reordena los candidatos\\
8 & \cod{combination} & Cómo se fusionan dos o más flujos en una respuesta\\
9 & \cod{routing}     & Qué flujo responde a qué panel\\
\bottomrule
\end{tabularx}
\captionof{table}{Los diez nodos. El orden es el del flujo de datos.}
\end{center}

\subsection{Cómo se cierra un nodo}

Las transiciones entre palabras son cuatro, y solo la última necesita algo que no
esté en las tablas de resultados:

\begin{center}
\begin{tabularx}{\linewidth}{@{}L{4.6cm} X@{}}
\toprule
transición & lo que hace falta\\
\midrule
\cod{blocked} $\rightarrow$ cualquiera & producir el artefacto que falta. Cada nodo bloqueado publica su propia precondición como un conteo leído del registro\\
\cod{inherited} $\rightarrow$ \cod{unpowered} & instanciar un segundo nivel. Se cuenta de las filas, nunca de una declaración\\
\cod{unpowered} $\rightarrow$ \cod{chosen} & puntuar el segundo nivel\\
\cod{chosen} $\rightarrow$ \cod{measured} & declarar un umbral que nombre un nivel real, y superarlo\\
\bottomrule
\end{tabularx}
\end{center}

\subsection{Por qué el grafo es la estructura y no las fases}

Una organización alternativa, y la que el proyecto usó antes, es por fases
sucesivas: fase A, fase B, y así. Tiene dos defectos frente al grafo.

El primero es que una fase no tiene tipo. Decir que la fase A está terminada no
dice si terminó porque se midió algo o porque se dejó de mirar. Un nodo con una
palabra de firmeza sí lo dice, y además no puede decirlo de más: la palabra se
calcula del registro y no se escribe a mano.

El segundo es que las fases sugieren un orden total donde solo hay un orden
parcial. El nodo \cod{routing} depende de que \cod{substrate} y \cod{retriever}
estén resueltos, pero no depende en absoluto de \cod{generator}. Una fase que
agrupe los tres obliga a esperar por el que no hace falta. El grafo permite que
seis nodos estén \cod{blocked} sin que eso impida avanzar en los otros cuatro,
que es exactamente la situación actual y no es un problema.
